\scanpage{34}%
\mspage{\handwritten{2}}%
% archive sheet number: 33
phisme, un monomorphisme, une immersion, une immersion ouverte, une immersion fermée, lisse, non ramifié ou étale (EGA IV 8.10.5, 11.2.6 et 17.7.6). Remarquons qu'il en est encore ainsi \deleted{\unsure{par|pour}} \uline{dans le cas où $\handwritten{\mathbb{P}}$ est la propriété d'être couvrant pour la topologie \handwritten{\unsure{/|f}}fidèlement plate localement de présentation finie} (notée fppf); en effet, étant donnés deux $S$-préschémas de présentation finie $X$ et $Y$, et un \deleted{morphisme}\inserted{$S$-morphisme} $u:X\handwritten{\longrightarrow}Y$, il résulte de (Exp. IV 6.3.1 (i))\deleted{,} que, pour que $u$ soit couvrant pour la topologie fppf, il faut et il suffit qu'il existe un $S$-préschéma $Z$ et un $S$-morphisme $v:Z\handwritten{\longrightarrow}Y$ fidèlement plat et de présentation finie qui se factorise à travers $u$.

Le but de ce numéro est \deleted{d'étendre}\inserted{\handwritten{de donner des variantes de}} ce genre de résultats \deleted{à}\inserted{\handwritten{pour}} la catégorie des $S$-groupes de présentation finie, \deleted{et à} celle des $S$-préschémas à groupe d'opérateurs, et \deleted{à}\inserted{\handwritten{pour}} certaines propriétés pour des monomorphismes de groupes (être invariant, central, à faisceau quotient inversible, etc ...).

Les deux résultats préliminaires de ce type sont les suivants :

\uline{10.\handwritten{2}.} \deleted{\unsure{a)}} \uline{Dans la situation rappelée au début de (10.0), soient \deleted{\unsure{pour}} $j\handwritten{\in}I$, $G_j$ et $H_j$ deux $S_j$-groupes de présentation finie; posons, pour tout $i\handwritten{\in}I$ tel que $j\handwritten{\leqq}i$, $G_i=G_j\handwritten{\times}_{S_j}S_i$, $G=G_j\handwritten{\times}_{S_j}S$, et définissons de même $\deleted{\unsure{G}}$ $H_i$ et $H$. Alors l'application canonique de $\underset{i\handwritten{\geqslant}j}{\underset{\handwritten{\longrightarrow}}{\lim}}\operatorname{Hom}_{\unsure{S|\mathbb{S}}_i\text{-gr.}}(G_i,H_i)$ dans $\operatorname{Hom}_{S\text{-gr.}}(G,H)$ est bijective.}

\uline{10.\handwritten{3}.} \deleted{\unsure{b)}} \uline{Dans la situation rappelée au début de (10.0), soit $G$ un $S$-groupe de présentation finie; alors il existe un indice $j\handwritten{\in}I$, \deleted{et} un $S_j$-groupe de présentation \deleted{f} finie $G_j$, et un \deleted{\illegible{1}} $S$-isomorphisme de groupes de $G$ sur $G_j\handwritten{\times}_{S_j}S$.}

Les assertions (10.\handwritten{\unsure{2|1}}) et (10.\handwritten{3}) sont des conséquences faciles de \deleted{(EGA IV 8.8.2 et 8.8.3)}\inserted{\handwritten{(10.0) et (10.1),}} compte tenu de l'interprétation de la structure de $S$-groupe donnée en (EGA $0_{\mathrm{III}}$ 8.2.5 et 8.2.6).

% Cancelled text visible through the pasted lower strip. The last line is
% read after vertical reflection; the other lines remain in the scan orientation.
\deleted{\unsure{10.4. Soit} $\unsure{\mathbb{P}}$ \unsure{l'une des propriétés} \illegible{5} \unsure{de préschémas en groupes} \illegible{1}}

\deleted{\unsure{(ii)|(i)} être un monomorphisme, \unsure{(i.e.} $\unsure{G}$ \illegible{7} $\unsure{H}$\unsure{)}}

\deleted{\unsure{(iii)|(ii)} être une immersion \unsure{(i.e.} $\unsure{G}$ \unsure{est un sous-préschéma en groupes de} $\unsure{H}$\unsure{)}}

\deleted{\unsure{(iv)|(iii)} être une immersion fermée \unsure{(i.e.} $G$ est un sous-préschéma en groupes fermé}

